\section{Density splitting and the taming cone}\label{sec:density-splitting}

The estimates of the preceding section have two consequences.  The first
requires no definite pair and imposes no condition on the cohomology class
of the resulting form. It is the fixed-structure theorem of
\emph{Taming implies compatibility on four-manifolds}
\cite[Theorem~1.1]{OpenAITaming2026}. The second consequence keeps a
specified class available, provided that it pairs positively with all
irreducible curves. This prescribed-class statement will be used along
the pair deformations in \Cref{sec:regular-preparation,sec:endpoint}.

\begin{theorem}[From taming to compatibility]\label{thm:tamed-compatible}
Let \(X\) be a closed connected smooth four-manifold and let \(J\) be a
smooth almost complex structure, with its complex orientation.  If a
smooth symplectic form tames \(J\), then there is a smooth symplectic form
compatible with \(J\).
\end{theorem}

In the following statement, curve
classes are identified with their Poincar\'e dual classes in
\(H^2(X;\R)\).  An irreducible \(J\)-holomorphic curve means the image of a
nonconstant somewhere-injective \(J\)-holomorphic map from a closed
connected Riemann surface, endowed with its complex orientation.

\begin{theorem}[A curve criterion for taming classes]\label{thm:taming-cone}
Let \(X\) be a closed connected oriented smooth four-manifold with
\(b^+(X)=3\), and let \(J\) be a smooth almost complex structure inducing
the given orientation.  Suppose that smooth closed real two-forms \(A,C\)
span a pointwise definite plane and form a basis of the
\(J\)-anti-invariant two-forms at every point.  Set
\[
 V=[A]^\perp\cap[C]^\perp\subset H^2(X;\R).
\]
Suppose \(U\in V\) has positive square and contains a smooth form taming
\(J\).  If \(H\in V\) has positive square, belongs to the same component of
\(\{v\in V:v^2>0\}\) as \(U\), and satisfies
\[
 H\cdot d>0
\]
for every irreducible closed \(J\)-holomorphic curve of class \(d\), then
\(H\) contains a smooth symplectic form taming \(J\).
\end{theorem}

We first combine \Cref{prop:angular-energy} with the tangent-measure
analysis to prove closedness of density weights in
\Cref{prop:density-weights}. The general splitting argument and its
smooth density-weight conclusion are shared with
\cite[Theorem~5.1 and Section~6, ``Passing to the positive-density set'']{OpenAITaming2026};
we give the argument with the conventions used here. The diffuse pairing then proves
\Cref{thm:tamed-compatible} without integral-current regularity.
Finally, the integral threshold cycles of
\Cref{lem:density-integral-thresholds} give \Cref{thm:taming-cone};
integral-current regularity enters only this last step.

\subsection{Densities and almost-everywhere planar tangents}

We use the conventions and separator constructions of
\Cref{lem:current-separators,prop:angular-energy}.  Thus \(g\) is a smooth
\(J\)-Hermitian metric with fundamental form \(G\),
\[
 N(\xi)=\int \xi_y(\ell)\,d\lambda(y,\ell),\qquad
 \mu=(\operatorname{pr}_X)_*\lambda,
 \qquad T=N+Q,
\]
where \(\ell\) ranges over the unit positively oriented complex-line
bivectors.  The current \(T\) is closed, \(Q\) is anti-invariant and belongs
to \(L^2\), and
\begin{equation}\label{eq:density-basic-bounds}
 \mu(B_s(p))\le Cs^2,\qquad
 \mathcal A_s:=
 \iint_{\dist(y,z)<s}|\ell-\tau_{zy}m|^2\,
          d\lambda(y,\ell)d\lambda(z,m)\le Cs^4.
\end{equation}
The map \(\tau_{zy}\) denotes parallel transport from \(z\) to \(y\) along
the short geodesic.  Constants in this section depend on the fixed
geometry and the current bounds, and may increase from line to line.
All radii are smaller than a fixed fraction of the injectivity radius.

Under the metric identification of coefficients, closedness of a
two-current means coclosedness of its distributional two-form.  We write
\(\mathfrak c(S)=[*S]\) for its Poincar\'e dual cohomology class when \(S\)
is closed.  In particular,
\begin{equation}\label{eq:density-intersection-convention}
 \mathfrak c(S)\cdot[\alpha]=S(\alpha)
\end{equation}
for a smooth closed two-form \(\alpha\).  This convention distinguishes
the current from its Hodge-dual differential-form representative.

The radial comparison below follows the Lelong--Jensen method for
almost-complex currents
\cite[Proposition~1 and Corollary~2]{ElkhadhraMimouni2007}.
Here we apply closedness to \(N+Q\) and estimate the \(L^2\)
correction \(Q\) directly.

\begin{lemma}[Existence of the quadratic density]\label{lem:density-existence}
At every \(p\in X\), the limit
\begin{equation}\label{eq:density-definition}
 \vartheta(p)=\lim_{s\downarrow0}s^{-2}\mu(B_s(p))
\end{equation}
exists and is finite.  The function \(\vartheta\) is measurable and
bounded, and it is zero almost everywhere with respect to \(dV_g\).
\end{lemma}

\begin{proof}
Put \(\nu=\mu+|Q|\,dV_g\).  Cauchy--Schwarz and the four-dimensional
volume bound give
\[
 \int_{B_s(p)}|Q|\,dV_g
 \le \vol_g(B_s(p))^{1/2}\|Q\|_{L^2(B_s(p))}
 \le Cs^2,
\]
so \(\nu(B_s(p))\le Cs^2\), uniformly in \(p\).
Fix \(p\), identify a normal-coordinate ball with a Euclidean ball, and
freeze \(J_p\) and \(G_p\) as constant tensors \(J_0,G_0\).  If
\(t=|\zeta|=\dist(p,\exp_p\zeta)\), define
\[
 b_p(s)=s^{-2}T(\mathbf1_{B_s(p)}G_0).
 \]
These evaluations are well-defined because \(T\) has measure
coefficients.  Since \(G(\ell)=1\), \(G-G_0=O(t)\), and \(Q\) annihilates
the invariant form \(G\), we have
\begin{equation}\label{eq:density-frozen-ratio}
 b_p(s)=s^{-2}\mu(B_s(p))+O(s).
\end{equation}
Indeed, both error integrals are bounded by
\(Cs^{-2}s\nu(B_s(p))\).

Call a radius good if its sphere has zero \(\nu\)-measure.  All but
countably many radii are good.  For good \(0<r<s\), closedness implies
\begin{equation}\label{eq:density-frozen-monotonicity}
 b_p(s)-b_p(r)
 =\frac12 T\bigl(\mathbf1_{\{r<t<s\}}\,dd^c_{J_0}\log t\bigr),
 \qquad d^c_{J_0}u=-du\circ J_0.
\end{equation}
For completeness, replace \(\log t\) inside \(B_s\) by
\[
 \phi_s(t)=
 \begin{cases}
 t^2/(2s^2)+\log s-1/2,&t<s,\\
 \log t,&t\ge s.
 \end{cases}
\]
The difference \(\phi_s-\phi_r\) is compactly supported and \(C^{1,1}\).
Inside \(B_s\), the quadratic part has
\(dd^c_{J_0}(t^2/(2s^2))=2G_0/s^2\).  Expanding
\(T(dd^c_{J_0}(\phi_s-\phi_r))=0\) gives
\Cref{eq:density-frozen-monotonicity}.
This use of \(C^{1,1}\) tests is justified by smoothing: for fixed
positive \(r,s\), their Hessians are uniformly bounded, converge away
from the two spheres, and the spheres have zero \(\nu\)-measure.

The frozen form \(dd^c_{J_0}\log t\) is semipositive on \(J_0\)-complex
lines and has norm at most \(Ct^{-2}\).  A \(J\)-complex line at distance
\(t\) is \(O(t)\) from a \(J_0\)-complex line.  Its evaluation can
therefore be negative only by \(C/t\).  Likewise the projection of this
form onto the actual anti-invariant summand has norm at most \(C/t\),
because its frozen anti-invariant projection vanishes.  Thus
\[
 b_p(s)-b_p(r)\ge
 -C\int_{0<t<s}t^{-1}\,d\nu.
\]
Quadratic growth, or a sum over the annuli
\(2^{-j-1}s<t\le2^{-j}s\), bounds this integral by \(Cs\).
It follows that
\begin{equation}\label{eq:density-almost-monotonicity}
 b_p(r)\le b_p(s)+Cs\qquad(0<r<s,\ r,s\text{ good}).
\end{equation}
The bounded function \(b_p\) therefore has a limit along good radii.
To see this directly, choose good \(s_j\downarrow0\) realizing its
limit inferior and apply \Cref{eq:density-almost-monotonicity} to all
smaller good \(r\); the limit superior is at most the same value.
Equation~\eqref{eq:density-frozen-ratio} gives the same limit for the
mass ratio.  Squeezing an arbitrary radius between good radii with
ratios tending to one gives \Cref{eq:density-definition}.

Measurability follows by taking a fixed sequence of radii tending to
zero, and boundedness follows from \Cref{eq:density-basic-bounds}.
For the last assertion, work in a fixed relatively compact coordinate
chart, localizing the measures before extending them to Euclidean
space. Apply the differentiation theorem for Radon measures, with
smooth four-dimensional volume as denominator
\cite[Theorem~4.19]{Kinnunen2026}. At volume-almost every center, the
measure-to-volume ratios of sufficiently small Euclidean coordinate
balls are bounded. A geodesic ball of radius \(s\) is contained in a
coordinate ball of comparable radius, whose smooth volume is
\(O(s^4)\). Hence \(\mu(B_s(p))=O_p(s^4)\) at those centers, and
\(s^{-2}\mu(B_s(p))\to0\). A finite chart cover proves the assertion
on \(X\).
\end{proof}

Write \(N=M\mu\), where \(M\) is the measurable positive Hermitian
bivector of trace one obtained by disintegrating \(\lambda\).  This
means \(M(p)=\int\ell\,d\lambda_p(\ell)\), where \(\lambda_p\) is a
probability measure on complex lines for \(\mu\)-almost every \(p\).

\begin{lemma}[Classification of the relevant tangent measures]
\label{lem:density-planar-tangents}
At \(\mu\)-almost every point \(p\) with \(\vartheta(p)>0\), the matrix
\(M(p)\) is a single unit complex-line bivector \(\ell_p\), the
conditional measure \(\lambda_p\) is supported on that line, and
\[
 s^{-2}(\exp_p^{-1}/s)_*\mu
 \ \longrightarrow\
 \frac{\vartheta(p)}{\pi}\,
 \mathcal H^2\!\restriction P_p
\]
locally weakly on \(T_pX\), where \(P_p\) is the plane of \(\ell_p\).
The pushforward here is restricted to the normal-coordinate ball before
rescaling.
\end{lemma}

\begin{proof}
Choose a positive-density point where differentiation of the bounded
matrix \(M\) with respect to \(\mu\) holds in a fixed Euclidean
coordinate chart and smooth local trivialization
\cite[Theorem~4.33 and Corollary~4.34]{Kinnunen2026}. This excludes a
set of zero \(\mu\)-measure. No doubling hypothesis is needed: for
fixed \(R\), enclose \(B_{Rs}(p)\) in a coordinate ball \(E_s\) of
radius comparable to \(Rs\). Quadratic growth bounds
\(s^{-2}\mu(E_s)\), while differentiation makes the mean oscillation
of \(M\) on \(E_s\) tend to zero. Therefore
\[
 s^{-2}\int_{B_{Rs}(p)}|M(y)-M(p)|\,d\mu(y)\longrightarrow0.
\]
A smooth change of frame adds a vanishing \(O(s)\) error after this
normalization. The
quadratic mass bound gives subsequential locally weak limits
\(\gamma\) of the rescaled scalar measures on all of \(T_pX\simeq\R^4\).
The existence of the density, followed by approximation of balls from
inside and outside, gives
\begin{equation}\label{eq:density-tangent-ball-mass}
 \gamma(B_R)=\vartheta(p)R^2\qquad(R>0).
\end{equation}
In particular every centered sphere has zero \(\gamma\)-measure.
The rescaled \(Q\) mass on that ball satisfies the sharper estimate
\begin{equation}\label{eq:density-rescaled-correction}
 s^{-2}\int_{B_{Rs}(p)}|Q|\,dV_g
 \le CR^2\|Q\|_{L^2(B_{Rs}(p))}\longrightarrow0.
\end{equation}
Passing to the limit in closedness of the rescaled \(T\) therefore
shows that the constant-coefficient current \(M(p)\gamma\) is closed.

Identify \(M(p)\) with a constant skew-symmetric matrix.  The equation
\(\partial(M(p)\gamma)=0\) says that all distributional derivatives
of \(\gamma\) in the column space of that matrix vanish.  Equivalently,
\(\gamma\) is translation invariant along \(\im M(p)\).
The nonzero positive Hermitian matrix \(M(p)\) has real rank two or
four.  Rank four would make \(\gamma\) a nonzero multiple of Lebesgue
measure on \(\R^4\), which contradicts
\Cref{eq:density-tangent-ball-mass}.  Thus its rank is two.
Trace-one normalization then makes \(M(p)=\ell_p\), the unit bivector
of a complex line \(P_p\).  The rank-one positive Hermitian matrix is
an extreme point of the trace-one positive matrices: if its
quadratic form vanishes on a vector, positivity forces almost every
matrix in any convex decomposition to vanish there too.  Hence
\(\lambda_p\) is concentrated on \(\ell_p\).

A locally finite measure invariant under translations in \(P_p\)
has the form
\(\mathcal H^2\!\restriction P_p\otimes\sigma\), with \(\sigma\) a
locally finite measure on \(P_p^\perp\).  Equation
\eqref{eq:density-tangent-ball-mass} now becomes
\[
 \vartheta(p)=
 \pi\int_{P_p^\perp}
       \bigl(1-|z_\perp|^2/R^2\bigr)_+\,d\sigma(z_\perp)
 \qquad(R>0).
\]
Each integrand is nondecreasing in \(R\), and is strictly increasing
once \(R>|z_\perp|>0\).  Constancy for all \(R\) forces \(\sigma\)
to be supported at zero.  Its mass is \(\vartheta(p)/\pi\).
Every convergent subsequence has consequently the same limit, proving
the asserted convergence.
\end{proof}

\begin{corollary}[Integrated transverse estimate]
\label{cor:density-transverse}
With the plane projection taken in \(T_yX\), one has
\begin{equation}\label{eq:density-transverse}
 \mathcal T_s:=
 \int d\lambda(y,\ell)
 \int_{B_s(y)}
 \left|\operatorname{pr}_{\ell^\perp}
                \exp_y^{-1}(z)/s\right|^2\,d\mu(z)
 =o(s^2).
\end{equation}
\end{corollary}

\begin{proof}
At the positive-density points covered by
\Cref{lem:density-planar-tangents}, divide the inner integral by
\(s^2\) and pass to the planar tangent measure.  The integrand vanishes
on that plane, and the limiting sphere has zero measure.
At zero-density points the same normalized integral is bounded by
\(s^{-2}\mu(B_s(y))\), which tends to zero.  These two arguments apply
for \(\lambda\)-almost every \((y,\ell)\).  The quadratic mass bound
gives a uniform integrable bound, so dominated convergence proves
\Cref{eq:density-transverse}.
\end{proof}

\subsection{Closedness of density weights}

The tangent measures now show that displacement normal to the local
complex line has a vanishing averaged contribution. We combine this
with the angular estimate to control derivatives of a smoothed density.
That control will allow the density weights to preserve closedness.

\begin{proposition}[Closed density weights]\label{prop:density-weights}
For either separator construction of \Cref{lem:current-separators},
let \(\vartheta\) be the density in \Cref{lem:density-existence}.
If \(b\colon\R\to\R\) is smooth and bounded with \(b(0)=0\), then
the current \(b(\vartheta)N\) is closed.
The currents
\begin{equation}\label{eq:density-weighted-restrictions}
 N^+=\mathbf1_{\{\vartheta>0\}}N,
 \qquad
 I_a=\mathbf1_{\{\vartheta>a\}}\frac{\pi}{\vartheta}N
 \quad(a>0)
\end{equation}
are also closed.  In particular, if
\(N^d=\mathbf1_{\{\vartheta=0\}}N\), then \(N^d+Q\) is closed.
\end{proposition}

\begin{proof}
We first smooth the scalar density.  Fix a nonnegative function
\(h\in C_c^\infty((0,1))\), normalized by
\(\int_0^1h(u)\,du=1\).  For small \(s\), put
\begin{equation}\label{eq:density-smoothed}
 f_s(y)=s^{-2}\int_X
          h\bigl(\dist(y,z)^2/s^2\bigr)\,d\mu(z),
 \qquad M_s(z)=\mu(B_s(z)).
\end{equation}
The function \(f_s\) is smooth: the support lies inside the injectivity
radius, and the cutoff vanishes near its boundary.  The mass bound
makes \(f_s\) uniformly bounded.  Integration with respect to the
radial distribution of \(\mu\) gives
\[
 f_s(y)=-\int_0^1 h'(u)\,
             \frac{\mu(B_{s\sqrt u}(y))}{s^2}\,du
 \longrightarrow
 -\vartheta(y)\int_0^1u h'(u)\,du=\vartheta(y).
\]
There are no atoms by quadratic growth, so no endpoint term occurs
at \(u=0\).  Thus \(f_s\to0\) volume-almost-everywhere as well.

Fix a smooth one-form \(\alpha\). To prove that \(b(\vartheta)N\)
is closed, apply closedness of \(T=N+Q\) before passing to the
density limit:
\begin{equation}\label{eq:density-weight-product-rule}
 0=T\bigl(d(b(f_s)\alpha)\bigr)
   =T\bigl(b(f_s)d\alpha\bigr)
    +T\bigl(b'(f_s)\,df_s\wedge\alpha\bigr).
\end{equation}
The first term has the required limit. Indeed, dominated convergence
gives
\[
 b(f_s)N\longrightarrow b(\vartheta)N,\qquad
 b(f_s)Q\longrightarrow0
\]
as currents. The second convergence uses \(b(0)=0\), the
volume-almost-everywhere convergence of \(f_s\), and domination
by \(C|Q|\).

It remains to show that the derivative term in
\eqref{eq:density-weight-product-rule} vanishes. The derivative of
\(b\) is bounded on the common bounded range of the \(f_s\).
Write \(|df_s|_\ell\) for the norm of the restriction to the plane
of \(\ell\). The \(N\) contribution is bounded by a constant times
\(\|\alpha\|_{C^0}\int |df_s|_\ell\,d\lambda\), and the \(Q\)
contribution is bounded by a constant times
\(\|\alpha\|_{C^0}\int |Q|\,|df_s|\,dV_g\).
We therefore need exactly two estimates:
\begin{equation}\label{eq:density-two-gradient-estimates}
 \int_X |Q|\,|df_s|\,dV_g=o(1),
 \qquad
 \int |df_s|_\ell\,d\lambda(y,\ell)=o(1).
\end{equation}
\smallskip
\noindent\emph{The derivative against \(Q\).}
Differentiating the radial kernel gives
\(|df_s(z)|\le Cs^{-3}M_s(z)\).  Fubini and the volume bound give
\[
 \int_X M_s\,dV_g
 =\int_X\vol_g(B_s(y))\,d\mu(y)\le Cs^4.
\]
Moreover \(M_s/s^2\) is uniformly bounded and tends to zero
volume-almost-everywhere.  Therefore
\begin{align}
 s^{-3}\int_X |Q|M_s\,dV_g
 &\le
 \left(\int_X |Q|^2\frac{M_s}{s^2}\,dV_g\right)^{1/2}
 \left(s^{-4}\int_X M_s\,dV_g\right)^{1/2}\notag\\
 &=o(1).
 \label{eq:density-q-gradient}
\end{align}
The first factor tends to zero by dominated convergence against
\(|Q|^2dV_g\), and the second is bounded.  This proves the first
estimate in \Cref{eq:density-two-gradient-estimates}.

\smallskip
\noindent\emph{The tangential derivative.}
Closedness will convert tangential differentiation into normal
derivatives. The radial kernel then supplies a transverse displacement
factor, multiplied by the difference between nearby complex-line
directions. The transverse estimate of \Cref{cor:density-transverse}
and the angular-energy bound will control these two factors together.

Fix a center and line \((y,\ell)\).  Choose an orthonormal basis of
\(T_yX\) whose first two vectors span the oriented plane of \(\ell\),
and let \(\zeta_1,\ldots,\zeta_4\) be the corresponding normal
coordinates.  The estimates below are uniform in this choice, so a
measurable choice of such bases suffices when integrating in
\((y,\ell)\).  Write
\[
 k_s^y(z)=h(|\zeta|^2/s^2),\qquad
 q_i(z)=\partial_{\zeta_i}k_s^y(z),\qquad
 \omega_{12}=d\zeta_1\wedge d\zeta_2.
\]
First variation of squared distance yields
\begin{equation}\label{eq:density-first-variation}
 df_s(y)(e_i)=-s^{-2}\int q_i(z)\,d\mu(z),
 \qquad i=1,2.
\end{equation}
The sign comes from differentiating the center; the displayed
\(q_i\) differentiates the radial expression in its coordinate
variable.

We replace this scalar integral by \(T(q_i\omega_{12})\).
On \(B_s(y)\), normal coordinates and parallel transport differ by
\(O(s^2)\), and hence
\begin{align}
 \omega_{12}(m)
 &=\langle\ell,\tau_{zy}m\rangle+O(s^2),\notag\\
 |1-\omega_{12}(m)|
 &\le C\bigl(s^2+|\ell-\tau_{zy}m|^2\bigr).
 \label{eq:density-tangent-form-error}
\end{align}
The second identity uses that both bivectors have unit length.
Since \(|q_i|\le C/s\), the integrated replacement error is at most
\begin{align}
 &s^{-2}\int
 \left|\int q_i\,d\mu-T(q_i\omega_{12})\right|d\lambda(y,\ell)
 \notag\\
 &\hspace{1cm}\le
 Cs^{-3}\left(
 s^2\iint_{\dist(y,z)<s}d\mu(y)d\mu(z)+\mathcal A_s\right)
 +Cs^{-3}\int_X|Q|M_s\,dV_g
 =o(1).
 \label{eq:density-replacement-error}
\end{align}
Here the double mass is \(O(s^2)\), the angular integral is \(O(s^4)\),
and the last term tends to zero by \Cref{eq:density-q-gradient}.
The normal-coordinate term is consequently \(O(s)\).

For \(i=1\), closedness of \(T\) on \(d(k_s^y\,d\zeta_2)\) gives
\[
 T(q_1\omega_{12})
 =-\sum_{a=3}^4T(q_a\,d\zeta_a\wedge d\zeta_2).
\]
For \(i=2\), using \(d(k_s^y\,d\zeta_1)\) gives the analogous
formula, with an irrelevant overall sign.  These are legitimate
smooth compactly supported tests because the cutoff vanishes near
the boundary of the coordinate ball.
The mixed forms satisfy, for \(k=1,2\) and \(a=3,4\),
\begin{equation}\label{eq:density-mixed-form-error}
 |(d\zeta_a\wedge d\zeta_k)(m)|
 \le C\bigl(s^2+|\ell-\tau_{zy}m|\bigr).
\end{equation}
Their value on the transported center line is zero, accounting for
the linear angular bound.  The normal derivative of the radial
kernel has the sharper estimate
\begin{equation}\label{eq:density-normal-derivative}
 |q_a(z)|\le
 Cs^{-1}
 \left|\operatorname{pr}_{\ell^\perp}\zeta/s\right|.
\end{equation}
Replacing \(T\) by \(N+Q\) in these mixed terms, the \(Q\) contribution
again tends to zero by \Cref{eq:density-q-gradient}; the \(s^2\)
geometric error contributes \(O(s)\).  Cauchy--Schwarz for the two
remaining angular and transverse factors bounds the rest by
\begin{equation}\label{eq:density-angular-transverse-product}
 Cs^{-3}\mathcal A_s^{1/2}\mathcal T_s^{1/2}
 =s^{-3}O(s^2)o(s)=o(1).
\end{equation}
Combining
\Cref{eq:density-first-variation,eq:density-replacement-error,eq:density-angular-transverse-product}
proves the second estimate in
\Cref{eq:density-two-gradient-estimates}.

\smallskip
\noindent\emph{The closed weighted limit.}
The two derivative estimates now allow passage to the limit in
\eqref{eq:density-weight-product-rule}. Its derivative term tends
to zero, and its first term tends to
\(b(\vartheta)N(d\alpha)\). Hence this last expression vanishes
for every smooth one-form \(\alpha\), proving the first assertion.

The remaining weights are obtained in a separate limit, after this
argument has been completed for each fixed smooth \(b\).  Choose a
smooth nondecreasing function \(\chi\colon\R\to[0,1]\) which is zero
on \((-\infty,0]\) and one on \([1,\infty)\).
The smooth bounded functions \(\chi(nt)\) converge pointwise to
\(\mathbf1_{\{t>0\}}\) on \([0,\infty)\), and vanish at zero.
For fixed \(a>0\), the functions
\[
 b_{a,n}(t)=
 \begin{cases}
 \pi\,\chi(n(t-a))/t,&t>0,\\
 0,&t\le0
 \end{cases}
\]
are smooth, vanish near zero, and are uniformly bounded by \(\pi/a\).
They converge at every \(t\ge0\) to
\(\mathbf1_{\{t>a\}}\pi/t\), with value zero at \(t=a\).
Dominated convergence and closedness of each smooth-weight current
prove \Cref{eq:density-weighted-restrictions}.  No derivative bound
uniform in \(n\) is used: the smoothing limit was taken first for
each \(n\).  Finally \(N^d+Q=T-N^+\) is closed.
\end{proof}

\subsection{The diffuse pairing and compatible forms}

We can now subtract the closed positive-density part of the current.
On the remaining zero-density set, the negative kernel error tends to
zero. The resulting intersection pairing first proves compatibility
and then gives the nonnegative square needed for the class criterion.

\begin{lemma}[Vanishing of the diffuse kernel error]
\label{lem:density-diffuse-error}
Set
\[
 B_s(z)=s^{-2}\int_X(1+\dist(y,z)/s)^{-6}\,d\mu(y).
\]
These functions are uniformly bounded, and \(B_s(z)\to0\) whenever
\(\vartheta(z)=0\).  In particular,
\begin{equation}\label{eq:density-diffuse-error}
 \int_X B_s(z)\,d\mu^d(z)\longrightarrow0,\qquad
 \mu^d=\mathbf1_{\{\vartheta=0\}}\mu.
\end{equation}
\end{lemma}

\begin{proof}
For each dyadic shell, quadratic growth gives the summable bound
\[
 s^{-2}(1+2^j)^{-6}\mu(B_{2^{j+1}s}(z))
 \le C\,2^{-4j}.
\]
The same estimate holds once the radius exceeds a fixed small
radius, by increasing \(C\) and using the finite total mass.
This proves uniform boundedness.  If \(\vartheta(z)=0\), each fixed
shell contribution tends to zero, because
\(s^{-2}\mu(B_{2^{j+1}s}(z))\to0\).
The summable bound then proves pointwise convergence.  A second
application of dominated convergence proves
\Cref{eq:density-diffuse-error}.
\end{proof}

For closed currents \(S,S'\) under consideration, the smoothing
operator \(S_s=(1+s^2\Delta_g)^{-3}\) preserves their cohomology classes
and commutes with the Hodge star.  Hodge decomposition therefore gives
\begin{equation}\label{eq:density-smoothed-intersection}
 \langle S_sS,*S_sS'\rangle_{L^2}
 =\mathfrak c(S)\cdot\mathfrak c(S'),
\end{equation}
independently of \(s\).  Their Sobolev regularity after smoothing is
sufficient for this identity; it can also be obtained by an
additional heat regularization and passage to the limit.

\begin{proof}[Proof of \Cref{thm:tamed-compatible}]
Suppose that no compatible symplectic form exists.
The separator in \Cref{lem:current-separators} gives a nonzero positive
current \(N\), normalized to have trace mass one, and an
anti-invariant \(Q\) such that \(T=N+Q\) annihilates all smooth closed
two-forms.  In particular \(T\) is closed and
\(\mathfrak c(T)=0\).
The preceding results apply to this separator.

By \Cref{prop:density-weights}, \(T^d=N^d+Q\) is closed.  Equation
\eqref{eq:density-smoothed-intersection} gives
\begin{align}
 0
 &=\langle S_sT,*S_sT^d\rangle\notag\\
 &=\langle S_sN,*S_sN^d\rangle
   +\langle S_sN,S_sQ\rangle
   +\langle S_sQ,S_sN^d\rangle
   +\|S_sQ\|_2^2 .
 \label{eq:density-compatibility-pairing}
\end{align}
The self-duality of \(Q\) accounts for the signs in this expansion.
Both mixed terms tend to zero by \Cref{prop:angular-energy}.
The positive-current lower bound of
\Cref{lem:current-positive-pairing}, applied with second submeasure
\(\lambda^d=\mathbf1_{\{\vartheta=0\}}\lambda\), gives
\[
 \langle S_sN,*S_sN^d\rangle
 \ge -C\int_XB_s\,d\mu^d=o(1).
\]
Since \(S_sQ\to Q\) in \(L^2\), taking a limit inferior in
\Cref{eq:density-compatibility-pairing} forces \(\|Q\|_2^2=0\).
Thus \(T=N\) is a nonzero positive current annihilating all closed
forms.

Let \(\eta\) be a taming symplectic form.  Compactness of the bundle
of unit complex lines and strict taming give
\(\eta_y(\ell)\ge c_\eta>0\).  Consequently
\[
 N(\eta)\ge c_\eta\mu(X)>0,
\]
contradicting its annihilation property.
This proves compatible-form existence.  The proof has used neither
rectifiability nor a theorem representing integral cycles by curves.
\end{proof}

\begin{lemma}[Nonnegative diffuse square]\label{lem:density-diffuse-square}
If \(Q=0\), then \(N^d\) is closed and
\(\mathfrak c(N^d)^2\ge0\).
\end{lemma}

\begin{proof}
Closedness follows from \Cref{prop:density-weights}.  By
\Cref{eq:density-smoothed-intersection,lem:current-positive-pairing},
\[
 \mathfrak c(N^d)^2
 =\langle S_sN^d,*S_sN^d\rangle
 \ge -C\int_X B_s\,d\mu^d.
\]
The right side tends to zero by
\Cref{lem:density-diffuse-error}.
\end{proof}

\subsection{Integral thresholds and positivity of the class}

To use the hypothesis on curve classes, we must turn the
positive-density current into integral cycles. The factor
\(\pi/\vartheta\) in \(I_a\) is chosen to make their multiplicity
one. Rectifiability and closedness will be checked before applying
the regularity theorem that produces curves.

\begin{lemma}[The threshold cycles]\label{lem:density-integral-thresholds}
For every \(a>0\), the current \(I_a\) in
\Cref{eq:density-weighted-restrictions} is an integral two-cycle with
positive \(J\)-complex tangent orientation.  If \(J\) is tamed, it is a
finite sum of currents of integration over irreducible closed
\(J\)-holomorphic curves, with positive integer multiplicities.
Moreover the identity
\begin{equation}\label{eq:density-layer-cake}
 N^+=\frac1\pi\int_0^\infty I_a\,da
\end{equation}
holds as an absolutely convergent integral of currents.
\end{lemma}

\begin{proof}
Put \(E_+=\{\vartheta>0\}\) and \(\mu^+=\mathbf1_{E_+}\mu\).
At \(\mu^+\)-almost every \(p\), Euclidean coordinate-ball
differentiation of \(\mathbf1_{E_+}\) gives value one. The
coordinate-ball enclosure argument in
\Cref{lem:density-planar-tangents}, applied to this bounded function,
gives
\[
 s^{-2}(\mu-\mu^+)(B_{Rs}(p))\longrightarrow0
 \qquad(R<\infty\text{ fixed}).
\]
Thus \(\mu^+\) has the same positive finite two-density and the same
planar tangent measure as \(\mu\) at almost every such point.
Under a smooth coordinate change its limit is the
linear image of the same plane measure.  Thus, in any fixed coordinate
chart, the Euclidean two-density exists and is positive and finite
almost everywhere for the pushed-forward measure: every Euclidean
sphere has zero mass for the limiting plane measure. Preiss's
rectifiability theorem therefore makes \(\mu^+\) a two-rectifiable measure
\cite{Preiss1987}; see also the theorem's formulation in
\cite[p.~771, opening paragraph]{KenigPreissToro2009}.
On a countably rectifiable carrier \(E\) it has
the representation
\begin{equation}\label{eq:density-rectifiable-measure}
 \mu^+=\frac{\vartheta}{\pi}\,
                 \mathcal H_g^2\!\restriction E.
\end{equation}
Here \(\pi\) is the area of the Euclidean unit two-disk.  The density
identification follows from differentiation on a rectifiable
carrier.  Its approximate tangent plane agrees almost everywhere
with the unique plane in \Cref{lem:density-planar-tangents}.
Thus \(N^+\) is integration over \(E\) with that complex tangent
orientation and multiplicity \(\vartheta/\pi\).

Multiplying by the weight defining \(I_a\) cancels this multiplicity:
\(I_a\) is integration with integer multiplicity one over
\(E\cap\{\vartheta>a\}\).  Its mass is at most
\(\pi\mu(X)/a\), and its boundary vanishes by
\Cref{prop:density-weights}.  It is consequently an integral current:
it is integer rectifiable, has finite mass, and has the zero integral
current as its boundary.

Suppose now that \(J\) is tamed.  By the already proved
\Cref{thm:tamed-compatible}, a compatible symplectic form exists.
In particular, the local compatibility hypothesis of the
Rivi\`ere--Tian regularity theorem holds.  That theorem applies to
integer rectifiable almost complex cycles and represents them by
\(J\)-holomorphic curves, with only isolated singularities
\cite[Theorem~I.1]{RiviereTian2009}.
Integral multiplicity and complex orientation do not depend on the
choice of Hermitian metric, so we may use the metric of this
compatible symplectic form for that application.
Normalize the resulting curve by separating its local branches at
the isolated singularities; the local completion is part of the
curve representation \cite[Section~VIII]{RiviereTian2009}.
Its normalization components are closed. Compactness and local
finiteness give only finitely many global components; equivalently,
small-ball monotonicity gives a uniform positive lower area for each
nonconstant closed component, whereas the total mass is finite.
Factoring each component through its somewhere-injective image gives
the asserted finite sum of irreducible closed curves with positive
integer multiplicities.

Finally, pointwise on \(\{\vartheta>0\}\),
\[
 \int_0^\infty
       \mathbf1_{\{\vartheta>a\}}\frac{\pi}{\vartheta}\,da=\pi.
\]
For a smooth two-form \(\alpha\), Fubini is justified by
\[
 \int_0^\infty |I_a(\alpha)|\,da
 \le \pi\|\alpha\|_{C^0}\mu^+(X).
\]
This proves \Cref{eq:density-layer-cake} with its stated normalization.
\end{proof}

\begin{proof}[Proof of \Cref{thm:taming-cone}]
Assume that the class \(H\) has no taming representative.
The class separator of \Cref{lem:current-separators} gives a nonzero
closed positive current \(N\), with \(Q=0\), satisfying
\begin{equation}\label{eq:density-cone-separator}
 H\cdot\mathfrak c(N)\le0.
\end{equation}
Its positive-density and diffuse parts are separately closed.
The anti-invariance of \(A,C\) implies
\[
 \mathfrak c(N^d)\cdot[A]=N^d(A)=0,\qquad
 \mathfrak c(N^d)\cdot[C]=N^d(C)=0.
\]
Thus \(w=\mathfrak c(N^d)\) lies in \(V\), and
\(w^2\ge0\) by \Cref{lem:density-diffuse-square}.

If \(N^d\ne0\), choose a taming form \(\eta\) in \(U\).  Strict
taming on the compact bundle of unit complex lines gives
\[
 U\cdot w=N^d(\eta)>0.
\]
In particular \(w\ne0\).  Since \([A],[C]\) span a positive
two-plane and \(b^+=3\), the intersection form on \(V\) has signature
\((1,b^-)\).  The conditions \(w^2\ge0\) and \(U\cdot w>0\) put
\(w\) in the nonzero closed time cone determined by \(U\).
Every timelike \(H\) in that component has \(H\cdot w>0\).
Explicitly, take a unit timelike vector \(u\) in the direction of
\(U\), and write \(H=h_0u+h_-\), \(w=w_0u+w_-\), with orthogonal
negative components.  Then
\[
 h_0>|h_-|,\qquad w_0\ge|w_-|,\qquad w_0>0,
\]
so
\[
 H\cdot w\ge h_0w_0-|h_-|\,|w_-|>0.
\]
If \(N^d=0\), its pairing is instead zero.  This distinguishes the
nonzero diffuse case required for strict positivity.

For every \(a>0\), \Cref{lem:density-integral-thresholds} expresses
\(I_a\) as a positive integral sum of irreducible closed
\(J\)-holomorphic curves.  The hypothesis on \(H\) therefore gives
\[
 H\cdot\mathfrak c(I_a)\ge0,
\]
with strict inequality whenever \(I_a\ne0\).
If \(N^+\ne0\), the layer identity has positive total mass, so the
set of \(a\) for which \(I_a\ne0\) has positive Lebesgue measure.
Pairing \Cref{eq:density-layer-cake} with any smooth representative
of \(H\) then gives \(H\cdot\mathfrak c(N^+)>0\).
The pairing is zero when \(N^+=0\).
At least one of \(N^d,N^+\) is nonzero, and hence
\[
 H\cdot\mathfrak c(N)
 =H\cdot\mathfrak c(N^d)+H\cdot\mathfrak c(N^+)>0.
\]
This contradicts \Cref{eq:density-cone-separator} and proves the
theorem.  A smooth closed taming representative is symplectic by
pointwise nondegeneracy.
\end{proof}
